% !PRINTED 775
\weberpage{775}
From this follows the inadmissibility of our assumption that $f_n(x)$, in accordance with formula (11), is decomposable into two factors in the field $R(\alpha)$.

\begin{center}
\rule{4cm}{0.4pt}
\end{center}

With this it can now be proved anew and in general that the equation whose roots are all the primitive $m$th roots of unity, and whose degree is $\varphi(m)$, is irreducible even when $m$ is an arbitrary composite number. This is equivalent to the following theorem:

\gesperrt{If an integral function $\varphi(x)$ with rational coefficients vanishes when a primitive $m$th root of unity is put for $x$, then $\varphi(x)$ vanishes also when any other primitive $m$th root of unity is put for $x$.}

When $m$ is a power of a prime, the theorem has already been proved in what precedes. We therefore apply complete induction: we assume it established for $m$th roots of unity and derive it for $mn$th roots of unity, where $n = p^k$ and $p$ is a prime that does not divide $m$.

Every $mn$th root of unity $\beta$ can be represented in the form

\[
\beta = \alpha r
\]

where $\alpha$ is an $m$th and $r$ an $n$th root of unity. What is to be proved can be stated thus: if $\beta$ is a root of a rational equation $\Phi(x) = 1$, % [sic]
then all the other primitive roots of unity of degree $mn$ must be roots of this equation as well.

If, as above, $f_n(x) = 0$ is the equation whose roots are the quantities $r$, then, as we have proved, $f_n(x)$ is not decomposable into rational factors, nor into such as depend rationally on $\alpha$. The function $f_n(\alpha^{-1}x)$, which vanishes for $x = \beta$, can therefore likewise not be decomposed into rational factors, because from a decomposition of the form

\[
f_n(\alpha^{-1}x) = \varphi(x,\alpha)\ \psi(x,\alpha)
\]

there would follow:

\[
f_n(x) = \varphi(\alpha x,\alpha)\ \psi(\alpha x,\alpha),
\]

which cannot hold. Since, accordingly, $f_n(\alpha^{-1}x)$ has a root
%CONT>>
